\documentclass[conference]{IEEEtran}

\IEEEoverridecommandlockouts

\usepackage{amsmath,amssymb}
\usepackage{graphicx}
\usepackage{url}
\usepackage[hidelinks]{hyperref}
\usepackage{fontspec}
\usepackage{xeCJK}

\setmainfont{Noto Serif}
\setsansfont{Noto Sans}
\setmonofont{Noto Sans Mono}
\setCJKmainfont[AutoFakeSlant=0.2]{Noto Serif CJK SC}
\setCJKsansfont[AutoFakeSlant=0.2]{Noto Sans CJK SC}
\setCJKmonofont{Noto Sans Mono CJK SC}
\XeTeXlinebreaklocale "zh"
\XeTeXlinebreakskip = 0pt plus 1pt

\begin{document}

\title{功率状态条件化的气象-功率融合网络用于区域风电集群总功率预测}

\author{\IEEEauthorblockN{罗傲}
\IEEEauthorblockA{\textit{School of Information Science and Technology}\\
\textit{University of Science and Technology of China}\\
Hefei 230022, China}
\and
\IEEEauthorblockN{刘磊\textsuperscript{*}}
\IEEEauthorblockA{\textit{School of Information Science and Technology}\\
\textit{University of Science and Technology of China}\\
Hefei 230022, China\\
\textsuperscript{*}通讯作者：liulei13@ustc.edu.cn}
\and
\IEEEauthorblockN{赵宏伟}
\IEEEauthorblockA{\textit{School of Information Science and Technology}\\
\textit{University of Science and Technology of China}\\
Hefei 230022, China}
\and
\IEEEauthorblockN{陈秋菊}
\IEEEauthorblockA{\textit{School of Information Science and Technology}\\
\textit{University of Science and Technology of China}\\
Hefei 230022, China}
\and
\IEEEauthorblockN{李斌}
\IEEEauthorblockA{\textit{School of Information Science and Technology}\\
\textit{University of Science and Technology of China}\\
Hefei 230022, China}
}

\maketitle

\begin{abstract}
业务化区域风电集群总功率 24 h 预测需要同时利用未来气象预报和目标窗口前可用的场站运行状态。未来网格化气象预报覆盖整个预测窗口，是驱动风电变化的主要外生信息；但仅依赖气象预报容易忽略预测起点前的真实功率水平、短期惯性和多场站同步变化，导致预测窗口前段存在状态失配误差。本文提出一种功率状态条件化的气象-功率融合预测网络。该方法将后报气象和未来气象预报编码为可查询的气象记忆，将最近 16 步场站级功率通过多尺度图编码器映射为动态功率上下文 token，并在 memory 层面注入气象记忆，使 future queries 在状态条件化后的融合记忆上并行生成 24 步区域总功率预测。基于吉林区域 58 个风电场站和 Pangu-Weather 网格化气象预报的实验表明，本文方法在测试集上取得 \(0.0767\pm0.0006\) 的 MAE、\(0.1041\pm0.0009\) 的 RMSE 和 \(0.8330\pm0.0029\) 的 \(R^2\)，优于 Weather-only 和多种网格气象编码基线。逐预测步分析显示，功率状态条件化主要改善预测窗口前段，尤其在 H1 上显著降低误差；消融实验进一步验证了真实功率状态、多尺度图功率编码和气象记忆结构对融合预测的贡献。
\end{abstract}

\begin{IEEEkeywords}
区域风电预测, 集群总功率, 气象-功率融合, 功率状态条件化, 图神经网络, Pangu-Weather
\end{IEEEkeywords}

\section{引言}

风电已成为现代电力系统中的重要可再生能源 \cite{ref1,ref2}。然而，风电出力受天气过程、地形条件和场站运行状态共同影响，具有显著的间歇性和波动性，会增加调度计划、备用容量配置和实时平衡的难度。因此，准确的风电功率预测是降低风电不确定性影响、支撑电力系统安全经济运行的重要工具。

随着风电场以区域场群形式接入电网，预测任务不再局限于单个场站出力，而可以面向目标区域内多个风电场站的聚合出力。本文关注区域集群总功率预测，即在给定目标窗口前可用观测和未来气象预报的条件下，预测研究区域内风电场站的总出力序列。该任务中的近期功率观测不仅反映单个场站的短期惯性，还包含多场站共同运行状态和场站间耦合关系；未来网格化气象预报则提供覆盖研究区域的外部气象驱动。

已有研究已经从多个方向提升了场群级风电预测性能，包括 CNN-LSTM 混合结构、网格化表示、图神经网络和趋势感知多步模型等 \cite{ref1,ref2,ref3,ref4,ref5,ref6,ref7,ref8,ref9,ref10,ref11,ref12}。同时，网格化数值天气预报和深度气象预报模型为未来窗口提供了重要外生信息，已有融合框架通常通过特征拼接、并行编码器或注意力交互将气象信息与功率信息结合 \cite{ref13,ref14,ref15,ref16,ref17,ref18}，也有工作直接对网格化气象场进行编码 \cite{ref19,ref20,ref21}。本文使用盘古气象大模型（Pangu-Weather）\cite{ref22} 作为预报来源。尽管这些方法证明了气象预报和功率观测融合的有效性，但多数框架仍倾向于将二者作为相对对称的输入分支处理，没有充分体现二者在业务化日内预测中的不同角色。

本文的核心出发点是：在本文的业务化日内集群总功率预测设定中，未来气象预报和近期场站功率承担不同角色。未来气象预报覆盖整个预测窗口，是驱动未来风电变化的主要外生因素；近期场站功率只在目标窗口前可用，主要表征当前运行状态、场站间耦合和短期惯性。因此，二者不应仅通过对称拼接或并行编码进行融合，而应采用更明确的角色划分：以未来气象预报作为主要驱动，以近期功率动态作为状态条件，对未来气象驱动的读取过程进行条件化约束。

基于这一认识，本文提出一种功率条件化气象记忆网络。该模型首先将历史和未来网格化气象场编码为可被预测查询读取的气象记忆；随后，将近期场站级功率动态映射为样本特定的 dynamic power context tokens，并注入气象记忆空间；最后，由并行 future queries 读取功率条件化后的气象记忆，输出 24 步集群总功率预测结果。该设计实现了跨模态的非对称条件化融合，使模型在读取未来气象驱动时显式受到当前功率状态的约束。

本文在吉林区域风电场群数据集上验证所提出的框架。实验在严格日样本和滑窗协议下进行，预测目标为区域集群总功率。结果表明，所提出模型在整体精度上优于 Weather-only 和多种网格气象编码基线；逐步长分析进一步显示，近期功率上下文对预测窗口前段具有更明显的补充作用；消融实验验证了真实功率状态、功率动态编码和气象记忆结构的必要性。

本文的主要贡献如下：

\begin{enumerate}
\item 本文从业务化区域集群总功率预测的时间可用性出发，明确区分未来气象预报和近期场站功率的建模角色：前者作为预测窗口内的主要外生驱动，后者作为目标窗口前的当前运行状态和多场站耦合信息。该视角为后续的非对称跨模态融合提供了建模依据。
\item 本文提出了一种功率条件化气象记忆网络，将历史和未来网格化气象场编码为可被预测查询读取的气象记忆，并将近期场站级功率动态转换为样本特定的 dynamic power context tokens 注入气象记忆空间，实现面向未来气象驱动的状态条件化读取。
\item 本文在吉林区域风电场群集群总功率预测任务上进行了系统验证。实验在严格时间划分和一致评价设置下比较 Weather-only、CNN、3D CNN、ViT、Swin 和 Video Transformer 等基线，并通过消融实验、误差分布和逐预测步分析验证了真实功率状态、功率动态编码和气象记忆结构的有效性。
\end{enumerate}

\section{相关工作}

\subsection{区域风电功率预测与时空建模}

风电功率预测可以按照空间尺度划分为单场站预测、场群预测和区域预测。对于电力系统调度而言，目标区域内多个风电场站的聚合出力通常比单个场站出力更直接对应功率平衡需求。区域或场群尺度的预测不仅需要描述功率序列的时间变化，还需要利用不同场站之间的空间相关性以及区域天气过程的连续传播特征。

已有研究从多种时空建模角度提升了风电预测精度。早期深度学习方法利用 CNN-LSTM 等结构联合建模局部空间特征和时间依赖 \cite{ref3}，也有工作将风机或场站嵌入规则网格，使卷积网络能够处理空间分布信息 \cite{ref4}。近年来，图神经网络被广泛用于场群级风电预测，因为风电场站可以自然表示为图节点，节点之间的边可由地理距离、气象相关性或功率相关性构造。相关方法包括时空图神经网络与模式识别结合 \cite{ref5}、全局信息自适应感知图卷积 \cite{ref6}、大规模场群动态图建模 \cite{ref7}、可学习图结构和动态误差修正 \cite{ref8}、趋势切换机制 \cite{ref9}、双融合时空图模型 \cite{ref10}、趋势分类与空间信息集成 \cite{ref11}，以及区域短期预测中的层级图神经网络 \cite{ref2} 和持续学习时空图框架 \cite{ref12}。

这些研究证明了多场站时空相关性建模对区域或场群预测的重要性。然而，现有方法的重点通常放在图结构构造、历史功率动态或场站间空间依赖上；即使引入气象信息，也常被作为节点特征或辅助输入使用。对于业务化多步预测中未来气象预报和近期场站功率的不同角色，即未来气象作为预测窗口内主要外生驱动、近期功率作为目标窗口前状态条件，已有工作讨论仍不充分。

\subsection{气象预报驱动的风电功率预测}

在日内多步风电预测中，未来气象预报是不可替代的外生输入。传统数值天气预报（Numerical Weather Prediction, NWP）和深度气象预报模型虽然来源不同，但都可以提供预测窗口内的风场和背景气象信息。本文使用 Pangu-Weather \cite{ref22} 生成的网格化气象预报作为输入，其建模作用与 NWP-功率融合研究中的未来气象驱动一致。

已有方法通常将 NWP 与历史功率作为两类输入，通过特征拼接、并行编码器、多模态注意力或多任务学习进行融合。例如，M2TNet \cite{ref13} 和 M2STAN \cite{ref14} 将 NWP 与功率信息建模为多模态多任务信息流；围绕 NWP 本身的误差特征，已有研究讨论了 NWP 预报风速与实测风速之间的时间滞后 \cite{ref15}、NWP 风速误差容忍性 \cite{ref16}、多源时空特征融合 \cite{ref17}，以及校正 NWP 与空间相关性建模的结合 \cite{ref18}。除节点或序列形式的 NWP 特征外，也有研究直接编码网格化气象场，例如使用 3D CNN 提取 NWP 数据立方体特征 \cite{ref19}、采用自动化深度学习框架进行区域风电预测 \cite{ref20}，或将深度 CNN 与概率预测框架结合用于区域概率风电预测 \cite{ref21}。近期工作还探索了大语言模型在多模态风电预测中的应用 \cite{ref23}。

上述研究表明，未来气象预报对区域风电预测具有关键作用，也说明直接编码网格化气象场是有效的建模方向。但多数融合框架仍将气象信息和功率信息作为相对平行的输入分支，通过拼接、并行编码或注意力交互进行对称融合。本文关注的是另一种角色划分：未来网格化气象预报作为预测窗口内的主要驱动，近期场站功率作为当前状态条件，对未来气象记忆的读取过程进行约束。

\subsection{条件化记忆与非对称跨模态融合}

Transformer 及其变体为多步时间序列预测提供了重要基础。Informer \cite{ref24}、PatchTST \cite{ref27}、Crossformer \cite{ref28} 和 Non-stationary Transformers \cite{ref34} 分别从长序列注意力、时间 patch 表示、跨变量依赖和非平稳性角度改进序列预测。风电预测中也已有多种 Transformer 或图 Transformer 结构用于建模时空耦合关系 \cite{ref37,ref38,ref39,ref40,ref41,ref42,ref43}。这些方法说明，注意力机制和并行多步解码适合处理复杂的风电时序预测问题。

与本文解码器设计更直接相关的是 query-based decoding 和可被注意力读取的记忆结构。DETR 使用 learnable queries 通过 cross-attention 读取编码特征并生成结构化输出 \cite{ref35}；Perceiver 和 Perceiver IO 通过 latent array 与 query-based decoding 处理大规模结构化输入输出 \cite{ref65,ref36}。更早的 Memory Networks 和 End-to-End Memory Networks 也说明，可读写或可注意力访问的记忆表示能够作为模型推理的信息载体 \cite{ref60,ref61}。在多模态系统中，Flamingo 使用视觉 tokens 注入语言模型，FiLM 则通过条件化调制改变特征表示 \cite{ref66,ref70}。这些工作为“用条件信息影响主模态读取过程”提供了形式参照。

需要强调的是，本文并不进行 prompt learning、语言模型适配或外部检索增强。本文中的 dynamic power context tokens 由当前样本的场站级功率动态直接生成，并被注入气象记忆空间；future queries 读取的是由功率状态条件化后的气象记忆。与风电预测中常见的特征级注意力或双注意力结构 \cite{ref72,ref73} 不同，本文的关键在于将近期功率作为状态条件前置到气象记忆中，使预测查询在读取未来气象驱动时显式受到当前多场站运行状态的约束。因此，本文方法更准确地表述为面向区域风电预测的功率条件化气象记忆，而不是通用 prompt 方法或对称多模态融合。

\newcommand{\methodfigure}[4]{%
\begin{figure*}[!t]
\centering
\includegraphics[width=#1]{#2}
\par\smallskip
\textbf{#3} #4
\end{figure*}
}

\section{方法}

\subsection{问题定义与业务时间对齐}

本文关注区域风电功率预测，预测目标为集群总功率，而不是单个场站级输出。设预测日为 \(D\)，\(t_0\) 为北京时间 \(D\) 日 00:00，目标窗口内的区域聚合功率序列记为 \(\mathbf{y}\in\mathbb{R}^{L_o}\)。本文中，近期功率上下文长度为 \(L_p=16\)，后报气象（Hindcast）和未来预报气象长度分别为 \(L_h=16\) 和 \(L_f=24\)，输出预测步长为 \(L_o=24\)。因此，完整气象序列长度为 \(L_m=L_h+L_f\)。场站级功率作为条件信息输入模型，最终输出和评价对象是区域聚合功率序列。

本文的预测协议遵循目标窗口前真实可用的信息。在日样本评价协议中，对于预测日 \(D\)，最新可用的场站级功率上下文覆盖 \(D-1\) 日 08:00-23:00 BJT。Pangu-Weather 预报 \cite{ref22} 在 \(D-1\) 日 08:00 BJT 起报。因此，气象 lead 0-15 对应 \(D-1\) 日 08:00-23:00 BJT，lead 16-39 对应预测目标窗口，即 \(D\) 日 00:00-23:00 BJT。本文将同一次气象预报中时间上位于目标窗口之前的气象帧称为后报气象（Hindcast），将目标窗口内的气象帧称为未来预报气象。在滑窗训练协议中，目标起点可以逐小时移动，但保持相同的起报到目标窗口间隔和输入输出对齐方式。

气象输入包括后报气象段 \(\mathbf{X}^{h}\in\mathbb{R}^{L_h\times C\times H\times W}\) 和未来预报气象段 \(\mathbf{X}^{f}\in\mathbb{R}^{L_f\times C\times H\times W}\)，其中 \(h\) 表示 Hindcast，\(C\)、\(H\) 和 \(W\) 分别表示气象变量数量和空间网格尺寸。功率输入由目标窗口前可用的场站级功率记录构成，模型仅保留最近 \(L_p\) 步，记为 \(\mathbf{P}^{c}\in\mathbb{R}^{N\times L_p}\)。场站邻接矩阵 \(\mathbf{A}\in\mathbb{R}^{N\times N}\) 在样本构造阶段由近期场站功率序列计算得到：先估计场站间功率相关性，保留高于阈值的正相关连接，其余连接置零，并去除自连接。该矩阵用于描述当前样本中风电场站之间的空间耦合关系。

为便于阅读，本文方法中的主要符号说明如下：\(d\) 表示隐空间维度；\(K\) 表示 dynamic power context tokens 的数量；\(L_w\) 表示气象记忆 token 数量；\(L_d\) 表示 decoder block 数量；\(l_p\) 和 \(s_p\) 分别表示功率局部片段长度和步长；\(J\) 表示由此得到的功率片段数量。

\methodfigure{\textwidth}{../../paper_image_use/method/Fig01.png}{图 1.}{业务时间协议与数据集构造。该图展示目标窗口前的场站功率上下文、Pangu 起报时间、对应气象 lead time 和预测目标窗口。}


\subsection{整体框架}

本文提出一种用于业务化日内集群总功率预测的功率条件化气象记忆网络。模型设计具有非对称性：未来网格化气象预报提供预测窗口内的主要外生驱动，近期场站级功率提供当前运行状态、场站间耦合关系和短期惯性。本文不是在输出层前简单拼接气象和功率特征，而是将功率状态转换为动态上下文 token，并将其作为条件信息注入气象记忆，使预测查询在读取未来预报气象时受到当前场群状态约束。

如图 2 所示，整体框架包含三个部分。第一，Weather Encoder 对后报气象和未来预报气象进行解耦空间编码、时空混合和记忆构造，得到 weather memory tokens \(\mathbf{M}^{w}\)。第二，Graph Encoder 对最近 \(L_p\) 步场站级功率和邻接矩阵进行多尺度图编码，得到功率状态向量 \(\mathbf{z}^{p}\)，并进一步生成 dynamic power context tokens \(\mathbf{C}^{p}\)。第三，Query Decoder 将 \(\mathbf{C}^{p}\) 前置到 \(\mathbf{M}^{w}\) 形成 conditioned memory，并使用可学习 future queries 并行解码 \(L_o\) 步区域聚合功率预测。

\methodfigure{\textwidth}{../../paper_image_use/method/Fig02.png}{图 2.}{本文提出的三模块功率条件化气象记忆网络整体框架。}

\subsection{Weather Encoder}

Weather Encoder 的输入由后报气象和未来预报气象沿时间维组成：

\begin{equation}
\mathbf{X}=\operatorname{Concat}_{time}(\mathbf{X}^{h},\mathbf{X}^{f})\in\mathbb{R}^{L_m\times C\times H\times W}.
\end{equation}

如图 3 所示，编码器首先使用三个解耦分支处理不同类型的气象信息。风场分支选择风场相关变量，背景场分支选择背景场变量，时间变化分支对完整气象序列计算一阶差分：

\begin{equation}
\Delta\mathbf{X}_1=\mathbf{0}, \qquad
\Delta\mathbf{X}_t=\mathbf{X}_t-\mathbf{X}_{t-1}, \quad t=2,\ldots,L_m.
\end{equation}

三个分支均采用 frame-wise CNN，即对每个气象帧独立进行二维空间编码。每个分支由一个 \(3\times3\) 卷积、GELU 激活和 2 个 Res2D block 组成；每个 Res2D block 内含两层 \(3\times3\) 卷积、dropout 和残差连接。设三条分支输出分别为 \(\mathbf{B}^{w}\)、\(\mathbf{B}^{s}\) 和 \(\mathbf{B}^{\Delta}\)，其融合过程可写为：

\begin{equation}
\mathbf{F}_0=\Phi_f\left([\mathbf{B}^{w};\mathbf{B}^{s};\mathbf{B}^{\Delta}]\right),
\end{equation}

其中 \(\Phi_f\) 由两层 \(1\times1\) 卷积和中间 GELU 激活构成，用于将三分支特征映射到统一隐空间维度 \(d\)。随后，模型加入帧位置嵌入和阶段嵌入：

\begin{equation}
\mathbf{F}_1=\mathbf{F}_0+\mathbf{e}^{frame}+\mathbf{e}^{phase},
\end{equation}

其中阶段嵌入用于区分后报气象帧和未来预报气象帧。加入位置与阶段信息后，特征被送入时空编码主干。该主干由 2 个 ST-Res block、1 个空间下采样层、2 个 Spatial Mix block 和 2 个后置 ST-Res block 组成。ST-Res block 使用深度可分离三维卷积和 \(1\times1\times1\) 卷积进行时空混合，并通过残差连接保持特征稳定；Spatial Mix block 在每个时间帧上使用深度二维卷积和 \(1\times1\) 卷积进一步混合局部空间结构。记该时空编码过程为 \(\mathcal{E}_w(\cdot)\)：

\begin{equation}
\mathbf{Z}^{w}=\mathcal{E}_w(\mathbf{F}_1).
\end{equation}

最后，Memory Builder 将编码后的气象特征转化为固定长度记忆。具体而言，编码特征先经过 \(1\times1\times1\) 投影，再通过自适应池化得到 \(L_m\times H_m\times W_m\) 的记忆网格：

\begin{equation}
\hat{\mathbf{Z}}^{w}=\operatorname{Pool}_{L_m\times H_m\times W_m}\left(\operatorname{Conv}_{1\times1\times1}(\mathbf{Z}^{w})\right).
\end{equation}

模型为该记忆网格加入时间、行和列位置嵌入：

\begin{equation}
\mathbf{M}^{w}_{t,y,x}=\hat{\mathbf{Z}}^{w}_{t,y,x}+\mathbf{e}^{time}_{t}+\mathbf{e}^{row}_{y}+\mathbf{e}^{col}_{x}.
\end{equation}

展平时空网格后得到 weather memory：

\begin{equation}
\mathbf{M}^{w}\in\mathbb{R}^{L_w\times d}, \qquad L_w=L_mH_mW_m.
\end{equation}

因此，\(\mathbf{M}^{w}\) 同时保留了后报气象上下文、未来预报气象驱动、变量分支特征和局部气象变化信息。

\methodfigure{\textwidth}{../../paper_image_use/method/Fig03.png}{图 3.}{Weather Encoder 的解耦气象分支、时空编码主干和记忆构造过程。}

\subsection{Graph Encoder}

Graph Encoder 从最近 \(L_p\) 步场站级功率中提取当前运行状态，并将其转化为 dynamic power context tokens。输入为功率上下文 \(\mathbf{P}^{c}\in\mathbb{R}^{N\times L_p}\) 和样本级邻接矩阵 \(\mathbf{A}\)。邻接矩阵在数据构造阶段由近期场站功率相关性得到；在 GCN 内部，模型进一步加入自环并进行对称归一化：

\begin{equation}
\tilde{\mathbf{A}}=\mathbf{A}+\mathbf{I},\qquad
\hat{\mathbf{A}}=\tilde{\mathbf{D}}^{-1/2}\tilde{\mathbf{A}}\tilde{\mathbf{D}}^{-1/2}.
\end{equation}

图卷积层可写为：

\begin{equation}
\operatorname{GCN}(\mathbf{H})=\hat{\mathbf{A}}\mathbf{H}\mathbf{W}+\mathbf{b}.
\end{equation}

其中，\(\mathbf{H}\) 表示当前图卷积层的节点特征，\(\mathbf{W}\) 为可学习权重矩阵，\(\mathbf{b}\) 为偏置项。

如图 4 所示，Graph Encoder 同时构造 point、patch 和 global 三种时间尺度的图表示。Point scale 将每个时间步的场站功率作为节点特征，经过 point GCN 得到 \(L_p\) 个 point tokens：

\begin{equation}
\mathbf{r}^{pt}_{t}=\Pi_{pt}\left(\operatorname{vec}\left(G_{pt}(\mathbf{P}^{c}_{:,t};\mathbf{A})\right)\right), \quad t=1,\ldots,L_p.
\end{equation}

Patch scale 将功率序列划分为长度为 \(l_p\)、步长为 \(s_p\) 的局部时间片段。令 \(a_j=1+(j-1)s_p\)，\(\mathcal{I}_j=\{a_j,\ldots,a_j+l_p-1\}\)，\(J=\lfloor(L_p-l_p)/s_p\rfloor+1\)，则 patch tokens 为：

\begin{equation}
\mathbf{r}^{pat}_{j}=\Pi_{pat}\left(\operatorname{vec}\left(G_{pat}(\mathbf{P}^{c}_{:,\mathcal{I}_j};\mathbf{A})\right)\right), \quad j=1,\ldots,J.
\end{equation}

Global scale 直接使用完整 \(L_p\) 步功率窗口，得到一个 global token：

\begin{equation}
\mathbf{r}^{g}=\Pi_{g}\left(\operatorname{vec}\left(G_{g}(\mathbf{P}^{c};\mathbf{A})\right)\right).
\end{equation}

在当前实现中，point、patch 和 global 三个 GCN encoder 均由 3 层图卷积构成，随后分别通过线性投影映射到维度 \(d\)。加入位置编码后，三类 token 分别记为 \(\mathbf{H}^{pt}\)、\(\mathbf{H}^{pat}\) 和 \(\mathbf{H}^{g}\)。跨尺度交互由 4 个 cross-scale attention block 构成，每个 block 包含多头交叉注意力、FFN、残差连接和 LayerNorm。交互顺序先自底向上，再自顶向下：

\begin{equation}
\begin{aligned}
\mathbf{H}^{pat}_{bu} &= \operatorname{CSA}_{pt\rightarrow pat}(\mathbf{H}^{pat},\mathbf{H}^{pt}), \\
\mathbf{H}^{g}_{bu} &= \operatorname{CSA}_{pat\rightarrow g}(\mathbf{H}^{g},\mathbf{H}^{pat}_{bu}), \\
\mathbf{H}^{pat}_{td} &= \operatorname{CSA}_{g\rightarrow pat}(\mathbf{H}^{pat}_{bu},\mathbf{H}^{g}_{bu}), \\
\mathbf{H}^{pt}_{td} &= \operatorname{CSA}_{pat\rightarrow pt}(\mathbf{H}^{pt},\mathbf{H}^{pat}_{td}).
\end{aligned}
\end{equation}

其中 \(\operatorname{CSA}_{source\rightarrow target}(\cdot)\) 表示以 target tokens 为 query、以 source tokens 为 key 和 value 的跨尺度注意力。交互后的 global、patch 和 point tokens 被拼接为一个融合序列，并送入 1 层 Transformer fusion encoder。随后，模型展平融合序列并经过状态投影得到功率状态向量：

\begin{equation}
\begin{aligned}
\mathbf{z}^{p}
&=f_p\left(\operatorname{TransformerEnc}\left([\mathbf{H}^{g}_{bu};\mathbf{H}^{pat}_{td};\mathbf{H}^{pt}_{td}]\right)\right) \\
&\in\mathbb{R}^{d}.
\end{aligned}
\end{equation}

最后，Graph Encoder 使用两层 MLP adapter 将 \(\mathbf{z}^{p}\) 映射为 \(K\) 个 dynamic power context tokens：

\begin{equation}
\mathbf{C}^{p}=\operatorname{reshape}\left(a_{\theta}(\mathbf{z}^{p})\right)\in\mathbb{R}^{K\times d}.
\end{equation}

当前模型不叠加样本无关的静态 prompt basis，也不使用额外门控缩放；adapter 输出的 \(\mathbf{C}^{p}\) 直接作为样本特定的功率上下文 token。

\methodfigure{\textwidth}{../../paper_image_use/method/Fig04.png}{图 4.}{Graph Encoder 的多尺度 GCN、跨尺度注意力和动态功率上下文构造。}

\subsection{Query Decoder}

Query Decoder 负责从由气象记忆和功率上下文共同组成的 conditioned memory 中读取预测信息。首先，Graph Encoder 生成的 \(K\) 个 dynamic power context tokens 被前置到 Weather Encoder 输出的 weather memory：

\begin{equation}
\mathbf{M}^{pc}=[\mathbf{C}^{p};\mathbf{M}^{w}]\in\mathbb{R}^{(K+L_w)\times d}.
\end{equation}

这种注入方式使功率状态以 memory token 的形式参与后续 cross-attention，而不是仅在输出层前与气象特征拼接。解码器维护可学习 future queries：

\begin{equation}
\mathbf{Q}^{(0)}\in\mathbb{R}^{L_o\times d},
\end{equation}

其中每个 query 对应一个预测 horizon。当前实现使用 2 个 QueryDecoderBlock。第 \(\ell\) 个 block 依次执行 query self-attention、query-to-memory cross-attention 和 FFN：

\begin{equation}
\begin{aligned}
\tilde{\mathbf{Q}}^{(\ell)} &= \operatorname{LN}\left(\mathbf{Q}^{(\ell-1)}+\operatorname{MHA}(\mathbf{Q}^{(\ell-1)},\mathbf{Q}^{(\ell-1)},\mathbf{Q}^{(\ell-1)})\right), \\
\bar{\mathbf{Q}}^{(\ell)} &= \operatorname{LN}\left(\tilde{\mathbf{Q}}^{(\ell)}+\operatorname{MHA}(\tilde{\mathbf{Q}}^{(\ell)},\mathbf{M}^{pc},\mathbf{M}^{pc})\right), \\
\mathbf{Q}^{(\ell)} &= \operatorname{LN}\left(\bar{\mathbf{Q}}^{(\ell)}+\operatorname{FFN}(\bar{\mathbf{Q}}^{(\ell)})\right).
\end{aligned}
\end{equation}

经过 \(L_d\) 个 decoder block 后，query 表示经过 LayerNorm 和预测头得到最终输出：

\begin{equation}
\hat{\mathbf{y}}=g_{\phi}\left(\operatorname{LN}(\mathbf{Q}^{(L_d)})\right)\in\mathbb{R}^{L_o}.
\end{equation}

预测头由线性层、GELU、dropout 和输出线性层组成。由于所有 future queries 并行读取同一个 conditioned memory，模型能够同时利用未来预报气象驱动和当前功率状态条件生成多步预测结果。

\subsection{训练目标}

主模型仅使用最终预测目标进行端到端训练。损失函数由平均绝对误差和均方根误差加权组成：

\begin{equation}
\mathcal{L}=\alpha\,\operatorname{MAE}(\hat{\mathbf{y}},\mathbf{y})
+(1-\alpha)\,\sqrt{\operatorname{MSE}(\hat{\mathbf{y}},\mathbf{y})+\varepsilon}.
\end{equation}

其中，\(\alpha\) 用于平衡两类误差，\(\varepsilon\) 是保证数值稳定性的小常数。当前模型不使用历史辅助重构损失、连续性约束损失或 start-constraint loss。这样可以使训练目标与最终集群总功率预测任务保持一致。

\newcommand{\casefigure}[4]{%
\begin{figure*}[!t]
\centering
\includegraphics[width=#1]{#2}
\par\smallskip
\textbf{#3} #4
\end{figure*}
}

\newcommand{\casetable}[3]{%
\begin{table*}[!t]
\centering
\textbf{#1} #2
\par\smallskip
#3
\end{table*}
}

\section{Case Study}

\subsection{原始数据}

本文采用吉林省区域风电场群数据构建案例研究数据集。该数据集由两部分组成：一部分为风电场站的实测功率数据，另一部分为覆盖研究区域的网格化气象预报数据。实测功率数据用于描述风电场群的历史运行状态，气象预报数据用于提供预测窗口内的外部气象驱动。与单场站预测不同，本文关注区域层面的集群总功率预测，预测目标为研究区域风电场群的区域集群总功率。

功率数据覆盖 2022--2023 年吉林省多个风电场站的运行记录。研究区域风电场群包含 58 个风电场站。所有功率数据均以北京时间为时间基准，并统一整理为整点时间序列。研究区域内的场站分布如图 5 所示。

\casefigure{0.54\textwidth}{../../paper_image_use/data/fig05_station_distribution.png}{图 5.}{吉林区域风电场站空间分布。}

气象数据采用 Pangu-Weather 生成的逐小时预报产品。本文截取吉林省及其周边区域的网格化气象场作为模型输入，使气象信息能够覆盖风电场群所在的主要空间范围。气象变量包括海平面气压、2\,m 气温、不同气压层上的位势高度、比湿和温度，以及 10\,m、1000\,hPa 和 925\,hPa 风场。为更直接表示风电相关风况，风矢量变量被转换为风速及方向分量，最终形成 17 个气象输入通道。本文后续将这些变量按风场相关变量和背景气象变量组织，其中风场信息是影响风电出力变化的关键外部因素。每个样本中的气象预报按照第 3.1 节定义的业务时间协议与功率观测和预测目标进行对齐。

训练集与评估集按年份隔离划分。训练集使用 2022 年数据，并采用滑动窗口方式构造训练样本，以提高训练样本规模并覆盖不同运行状态。验证集和测试集使用 2023 年数据，均采用严格日样本设置，即每天 00:00 BJT 构造一个 24 小时预测样本。具体划分时，每月前 10 天作为验证集，其余日期作为测试集，验证样本和测试样本在预测日上互不重叠。该划分方式避免了随机划分可能造成的时间信息泄漏，也使模型选择和最终测试保持独立。

\subsection{实验设置}

为验证所提出模型在吉林区域风电场群 24 步集群总功率预测任务中的有效性，本文在固定的数据划分上进行实验。训练集、验证集和测试集分别对应第 4.1 节中给出的时间划分。所有模型均在相同的测试集上进行评价，并重复运行 12 次，报告均值和标准差。

与风电功率预测研究中常用的评价方式一致，本文采用 MAE、RMSE 和 \(R^2\) 评价确定性预测误差。设 \(P_i\) 和 \(\hat{P}_i\) 分别表示第 \(i\) 个测试点的真实集群总功率和预测集群总功率，\(\bar{P}\) 表示测试集中真实功率的均值，\(n\) 表示参与评价的测试点总数。在整体 24 步评价中，\(n\) 为测试样本数与预测步长的乘积；在逐步长评价中，\(n\) 为该预测步长下的测试样本数。由于模型输出为归一化功率，令 \(S_y\) 为训练集目标功率归一化尺度，\(y_i=P_i/S_y\)，\(\hat{y}_i=\hat{P}_i/S_y\)。归一化误差指标定义为：

\begin{equation}
\mathrm{MAE}=\frac{1}{n}\sum_{i=1}^{n}|y_i-\hat{y}_i|,
\end{equation}

\begin{equation}
\mathrm{MSE}=\frac{1}{n}\sum_{i=1}^{n}(y_i-\hat{y}_i)^2,
\end{equation}

\begin{equation}
\mathrm{RMSE}=\sqrt{\frac{1}{n}\sum_{i=1}^{n}(y_i-\hat{y}_i)^2}.
\end{equation}

实际功率尺度下的误差定义为：

\begin{equation}
\mathrm{MAE(MW)}=\frac{1}{n}\sum_{i=1}^{n}|P_i-\hat{P}_i|,
\end{equation}

\begin{equation}
\mathrm{RMSE(MW)}=\sqrt{\frac{1}{n}\sum_{i=1}^{n}(P_i-\hat{P}_i)^2}.
\end{equation}

\(R^2\) 用于衡量预测结果对测试集功率变化的解释能力，定义为：

\begin{equation}
R^2=1-\frac{\sum_{i=1}^{n}(P_i-\hat{P}_i)^2}{\sum_{i=1}^{n}(P_i-\bar{P})^2}.
\end{equation}

除整体 24 步平均结果外，本文还进一步分析不同预测步长上的误差变化，以检验模型在预测窗口内部的稳定性。

\casetable{表 1.}{实验设置。}{\input{../../paper_table_use/result/table01_experimental_setup.tex}}

对比模型包括 Weather-only 结构以及多种网格气象编码基线，包括 CNN、3D CNN、ViT、Swin 和 Video Transformer。所有基线均使用相同的数据划分和评价指标。本文模型与 Weather-only 的差异在于，本文模型进一步引入近期功率上下文，并将其用于条件化气象记忆；因此二者的对比用于检验功率状态信息是否能为未来气象驱动预测提供补充。

\subsection{基准模型对比}

表 2 给出了不同模型在测试集上的整体预测结果。本文模型在所有主要指标上均取得最优结果，MAE 为 \(0.0767\pm0.0006\)，RMSE 为 \(0.1041\pm0.0009\)，\(R^2\) 为 \(0.8330\pm0.0029\)。在实际功率尺度下，本文模型的 MAE(MW) 和 RMSE(MW) 分别为 \(377.9\pm3.1\) MW 和 \(512.7\pm4.4\) MW。

\casetable{表 2.}{基准模型整体预测结果。}{\input{../../paper_table_use/result/table02_benchmark_results.tex}}

与 Weather-only 相比，本文模型的 MAE 从 0.0800 降至 0.0767，RMSE 从 0.1072 降至 0.1041，\(R^2\) 从 0.8230 提升至 0.8330。这说明在相同气象输入条件下，引入近期功率上下文能够降低整体预测误差。除 Weather-only 外，3D CNN 是表现最好的气象编码基线，其 MAE 和 RMSE 分别为 0.0845 和 0.1138，仍高于本文模型。CNN、ViT、Swin 和 Video Transformer 的误差进一步增大。

\casefigure{\textwidth}{../../paper_image_use/result/fig06_benchmark_comparison.png}{图 6.}{基准模型误差对比。}

图 6 以误差条形式展示了各模型的 MAE 和 RMSE。图中排序与表 2 的整体结果一致：本文模型误差最低，Weather-only 位列第二，其余气象编码基线误差较高。Weather-only 与本文模型的差距小于其他气象编码基线，表明气象预报本身已经提供了强预测信号；但在此基础上加入功率上下文后，模型仍能进一步改善测试集表现。

\subsection{消融实验}

表 3 给出了功率上下文、功率编码器和气象模块相关消融结果。所有消融模型均与本文完整模型在相同测试集上比较，表中 \(\Delta\)MAE 和 \(\Delta\)RMSE 表示相对于完整模型的误差增幅。

为便于解释表中 Variant 的含义，本文将消融设置说明如下。Weather-only 移除所有功率侧条件化，仅保留气象编码器和 Query Decoder；Random power 和 Zero power 保留功率条件通路，但在功率编码器接收最近 \(L_p\) 步功率后，分别将其替换为同形状随机值和零值；Basic GCN 将完整模型中的点级、片段级和全局多尺度图编码以及跨尺度注意力融合，替换为单个 GCN 对最近 \(L_p\) 步功率进行编码。

气象侧消融用于检验 Weather Encoder 的结构设计。No memory 关闭气象记忆的空间组织，将气象特征在空间维聚合为每个时刻一个 memory token，仅保留时间位置表示；No delta 移除气象一阶差分分支，仅使用风场分支和背景场分支；Single branch 将风场、背景场和差分三路解耦编码改为对全部气象通道使用单一分支；No synoptic 移除背景场分支，并在差分分支中屏蔽背景场通道；No wind 移除风场分支，并在差分分支中屏蔽风场通道。

\casetable{表 3.}{消融实验结果。}{\input{../../paper_table_use/result/table03_ablation_results.tex}}

首先，去除功率上下文后的 Weather-only 模型 MAE 增加 4.27\%，RMSE 增加 2.94\%。将功率输入替换为随机值或零值后，MAE 分别增加 4.72\% 和 4.70\%，与 Weather-only 的退化幅度接近。这说明模型性能提升来自真实功率状态信息，而不是额外分支或参数量本身。将功率编码器替换为 Basic GCN 后，MAE 增加 2.29\%，RMSE 增加 1.96\%，表明图结构功率建模是有效的，但完整的多尺度功率动态编码能够取得更低误差。

其次，气象模块消融显示，气象空间记忆和气象变化建模对预测结果有直接影响。关闭气象空间记忆后，MAE 和 RMSE 分别增加 5.33\% 和 5.01\%；去除气象差分分支后，MAE 和 RMSE 分别增加 4.62\% 和 4.83\%；将解耦气象分支合并为单分支后，MAE 和 RMSE 分别增加 8.88\% 和 8.76\%。这些结果说明，相应气象编码组件被移除后，测试集误差均出现上升。

最后，气象通道消融表明风场信息对测试集误差影响最大。屏蔽背景气象变量后，MAE 和 RMSE 分别增加 11.66\% 和 10.95\%；屏蔽风场相关变量后，MAE 和 RMSE 分别增加 97.79\% 和 96.58\%，\(R^2\) 降至 0.3538。这说明在本文数据集上，风场相关变量是不可缺少的气象输入。

\casefigure{\textwidth}{../../paper_image_use/result/fig07_ablation_delta_vs_main.png}{图 7.}{消融模型相对于完整模型的误差增幅。}

\subsection{误差分布分析}

图 8 给出了不同模型在测试样本上的绝对误差分布。本文模型的误差分布整体低于其他基线，说明其优势不仅体现在平均指标上，也体现在测试样本层面的误差分布上。与此同时，各模型均存在较长的高误差尾部，表明在部分复杂天气或功率波动样本上，24 步集群总功率预测仍具有较高难度。

\casefigure{\textwidth}{../../paper_image_use/result/fig08_error_distribution.png}{图 8.}{不同模型的绝对误差分布。}

\subsection{不同预测步长上的结果}

为进一步分析模型在 24 步预测窗口内的表现，表 4 比较了本文模型和 Weather-only 在 H1、H6、H12、H24 以及全预测窗口平均意义下的误差。图 9 展示了两个模型在所有预测步长上的 MAE、RMSE 以及 MAE 差值。

\casetable{表 4.}{不同预测步长上的误差对比。}{\input{../../paper_table_use/result/table04_horizon_summary.tex}}

在首个预测步 H1 上，本文模型的 MAE 为 0.0430，显著低于 Weather-only 的 0.0747；RMSE 也从 0.0956 降至 0.0586。该结果说明近期功率上下文对预测窗口开始阶段最为有效，因为此时当前运行状态和短期惯性对目标功率仍具有较强约束作用。

随着预测步长增加，本文模型与 Weather-only 之间的差距逐渐缩小。在 H6 和 H12 上，本文模型仍保持较小优势；到 H24 时，Weather-only 的 MAE 和 RMSE 略低于本文模型。该结果表明，功率上下文带来的误差降低主要集中在预测窗口前段。尽管如此，在 1--24 步平均意义下，本文模型仍保持更低的 MAE 和 RMSE。

\casefigure{\textwidth}{../../paper_image_use/result/fig09_ours_weather_horizon_residual.png}{图 9.}{本文模型与 Weather-only 的逐步长误差及差值。}

\subsection{Discussion}

实验结果首先说明，区域集群总功率预测中未来气象预报仍是主导性信息来源。Weather-only 在所有基线中表现较强，表明 Pangu-Weather 网格化气象场已经能够刻画预测窗口内的主要外生驱动。然而，仅依赖气象场并不能完全描述目标窗口开始前的实际运行状态。本文模型在 Weather-only 的基础上进一步引入近期场站级功率，并将其作为动态上下文注入气象记忆，使预测过程不只是从未来气象场直接回归功率，而是在当前场群状态约束下读取未来气象信息。

逐步长结果揭示了这一条件化机制的作用边界。功率上下文带来的改善主要集中在预测窗口前段，尤其在 H1 上最为明显；随着预测步长增加，本文模型与 Weather-only 的差距逐渐缩小，H24 上 Weather-only 甚至略优。这一现象并不削弱功率上下文的意义，而是说明不同信息源在时间轴上的作用不同：目标窗口前最新功率更适合约束近端预测，因为其中包含当前出力水平、短期惯性和场站间同步变化；远端预测则更多受未来气象预报质量支配。因此，本文方法的核心不是用历史功率替代气象预报，而是用近期功率对气象驱动预测进行状态条件化。

消融实验进一步支持上述解释。随机功率和零功率输入接近 Weather-only，说明收益并非来自额外参数量或功率分支形式，而依赖真实功率状态本身。Basic GCN 的性能弱于完整模型，表明仅用单一时间尺度的图编码不足以充分刻画场站级功率动态；点级、片段级和全局尺度之间的交互有助于形成更有效的功率上下文。气象侧消融则表明，空间气象记忆、气象变化分支和变量解耦编码共同影响模型读取未来气象信息的能力。

气象通道消融还给出了较清晰的物理解释。屏蔽背景气象变量会造成误差上升，但屏蔽风场相关变量会导致性能大幅退化，说明本文任务中的主要可预测性仍来自未来风场。背景场变量的作用更接近对大尺度天气状态的补充约束，而不能替代与风电出力直接相关的风场信息。这一结果也表明，模型性能并非来自简单堆叠更多气象通道，而是依赖气象变量与风电出力之间的物理相关性。

从应用角度看，本文结果支持一种符合时间可用性约束的区域集群预测流程：在每天预测目标窗口开始前，使用最新可获得的场站级功率表征当前运行状态，同时使用 Pangu-Weather 预报提供未来气象驱动。该流程对预测窗口前段尤其有价值；对于更远预测步长，进一步结合气象预报不确定性有望提升模型对远端预测的适应能力。本文以吉林区域风电场群为案例开展系统验证，后续可在更多区域和不同装机结构的风电场群上检验模型迁移能力，并进一步面向极端爬坡、限电恢复和概率预测等调度相关场景扩展。

\section{结论}

本文面向区域风电场群 24 步集群总功率预测问题，研究了如何在严格时间可用性约束下联合利用近期场站级功率和未来网格化气象预报。本文的核心思路是对两类输入赋予不同角色：未来气象预报提供预测窗口内的主要外生驱动，近期功率记录表征目标窗口前的当前运行状态、场站间耦合和短期惯性。基于这一设定，本文提出功率条件化气象记忆网络，将近期功率动态编码为样本特定的 dynamic power context tokens，并注入气象记忆空间，使 future queries 能够在状态条件化后的气象记忆中读取预测信息。

吉林区域风电场群案例研究表明，所提出模型在整体 24 步预测任务上优于 Weather-only 以及多种网格气象编码基线，并在 MAE、RMSE 和 \(R^2\) 上取得更优结果。逐步长分析显示，近期功率上下文的贡献主要集中在预测窗口前段，这与其刻画当前运行状态和短期惯性的作用一致；远端 horizon 上模型表现则更多受未来气象预报质量影响。

消融实验进一步验证了模型设计的有效性。真实功率输入、多尺度图功率编码、气象空间记忆、气象变化分支和变量解耦编码均对预测性能产生贡献；气象通道消融也表明，风场相关变量是该任务中不可替代的关键输入，背景气象变量主要提供补充性约束。

总体而言，本文结果说明，在区域风电集群总功率预测中，将近期功率作为动态状态上下文、将未来气象作为主要驱动信息，是一种有效且符合业务时间约束的建模方式。未来工作可进一步围绕跨区域迁移、气象预报不确定性建模，以及极端爬坡、限电恢复和概率预测等调度相关场景展开。

\begin{thebibliography}{43}

\bibitem[1]{ref1}
Wang L, Gao J, Li Y, Wang D. A method for ultra-short-term wind power forecasting of large-scale wind farms based on adaptive spatiotemporal graph convolution. Renew Energy 2025;249:123321. \url{https://doi.org/10.1016/j.renene.2025.123321}

\bibitem[2]{ref2}
Li M, Yang M, Yu Y. Short-term regional wind power forecast method based on hierarchical graph neural network and extreme value theory. Energy Convers Manag 2025;341:120002. \url{https://doi.org/10.1016/j.enconman.2025.120002}

\bibitem[3]{ref3}
Liu H, Mi X, Li Y. Smart deep learning based wind speed prediction model using wavelet packet decomposition, convolutional neural network and convolutional long short term memory network. Energy Convers Manag 2018;166:120-31. \url{https://doi.org/10.1016/j.enconman.2018.04.021}

\bibitem[4]{ref4}
Yu R, Liu Z, Li X, Lu W, Ma D, Yu M, et al. Scene learning: deep convolutional networks for wind power prediction by embedding turbines into grid space. Appl Energy 2019;238:249-57. \url{https://doi.org/10.1016/j.apenergy.2019.01.010}

\bibitem[5]{ref5}
Liu X, Zhang Y, Zhen Z, Xu F, Wang F, Mi Z. Spatio-temporal graph neural network and pattern prediction based ultra-short-term power forecasting of wind farm cluster. IEEE Trans Ind Appl 2024;60:1794-803. \url{https://doi.org/10.1109/TIA.2023.3321267}

\bibitem[6]{ref6}
Yang M, Ju C, Huang Y, Guo Y, Jia M. Short-term power forecasting of wind farm cluster based on global information adaptive perceptual graph convolution network. IEEE Trans Sustain Energy 2024;15:2063-76. \url{https://doi.org/10.1109/TSTE.2024.3397877}

\bibitem[7]{ref7}
Yang M, Wang D, Zhang W, Yv X. A centralized power prediction method for large-scale wind power clusters based on dynamic graph neural network. Energy 2024;310:133210. \url{https://doi.org/10.1016/j.energy.2024.133210}

\bibitem[8]{ref8}
Yang M, Guo Y, Huang T, Fan F, Ma C, Fang G. Wind farm cluster power prediction based on graph deviation attention network with learnable graph structure and dynamic error correction during load peak and valley periods. Energy 2024;312:133645. \url{https://doi.org/10.1016/j.energy.2024.133645}

\bibitem[9]{ref9}
Yang M, Huang Y, Guo Y, Zhang W, Wang B. Ultra-short-term wind farm cluster power prediction based on FC-GCN and trend-aware switching mechanism. Energy 2024;290:130238. \url{https://doi.org/10.1016/j.energy.2024.130238}

\bibitem[10]{ref10}
Qu Z, Li J, Hou X, Gui J. A D-stacking dual-fusion, spatio-temporal graph deep neural network based on a multi-integrated overlay for short-term wind-farm cluster power multi-step prediction. Energy 2023;281:128289. \url{https://doi.org/10.1016/j.energy.2023.128289}

\bibitem[11]{ref11}
Yang M, Jiang Y, Xu C, Wang B, Wang Z, Su X. Day-ahead wind farm cluster power prediction based on trend categorization and spatial information integration model. Appl Energy 2025;388:125580. \url{https://doi.org/10.1016/j.apenergy.2025.125580}

\bibitem[12]{ref12}
Zhang Q, Mei K, Zhou Q, Ding X, Dong Z, Tian M, et al. A wind farm cluster power prediction method based on globally aware dynamic spatiotemporal graph continual learning framework. Energy 2026;348:140591. \url{https://doi.org/10.1016/j.energy.2026.140591}

\bibitem[13]{ref13}
Wang L, He Y, Liu X, Li L, Shao K. M2TNet: multi-modal multi-task Transformer network for ultra-short-term wind power multi-step forecasting. Energy Rep 2022;8:7628-42. \url{https://doi.org/10.1016/j.egyr.2022.05.290}

\bibitem[14]{ref14}
Wang L, He Y. M2STAN: multi-modal multi-task spatiotemporal attention network for multi-location ultra-short-term wind power multi-step predictions. Appl Energy 2022;324:119672. \url{https://doi.org/10.1016/j.apenergy.2022.119672}

\bibitem[15]{ref15}
Liu C, Zhang X, Mei S, Zhou Q, Fan H. Series-wise attention network for wind power forecasting considering temporal lag of numerical weather prediction. Appl Energy 2023;336:120815. \url{https://doi.org/10.1016/j.apenergy.2023.120815}

\bibitem[16]{ref16}
Yang M, Guo Y, Huang T, Zhang W. Power prediction considering NWP wind speed error tolerability: a strategy to improve the accuracy of short-term wind power prediction under wind speed offset scenarios. Appl Energy 2025;377:124720. \url{https://doi.org/10.1016/j.apenergy.2024.124720}

\bibitem[17]{ref17}
Yang M, Han C, Zhang W, Wang B. A short-term power prediction method for wind farm cluster based on the fusion of multi-source spatiotemporal feature information. Energy 2024;294:130770. \url{https://doi.org/10.1016/j.energy.2024.130770}

\bibitem[18]{ref18}
Hu S, Xiang Y, Zhang H, Xie S, Li J, Gu C, et al. Hybrid forecasting method for wind power integrating spatial correlation and corrected numerical weather prediction. Appl Energy 2021;293:116951. \url{https://doi.org/10.1016/j.apenergy.2021.116951}

\bibitem[19]{ref19}
Higashiyama K, Fujimoto Y, Hayashi Y. Feature extraction of NWP data for wind power forecasting using 3D-convolutional neural networks. Energy Procedia 2018;155:350-8. \url{https://doi.org/10.1016/j.egypro.2018.11.043}

\bibitem[20]{ref20}
Keisler J, Le Naour E. WindDragon: automated deep learning for regional wind power forecasting. Environ Data Sci 2025;4:e19. \url{https://doi.org/10.1017/eds.2025.10}

\bibitem[21]{ref21}
Jonkers J, Avendaño DN, Van Wallendael G, Van Hoecke S. A novel day-ahead regional and probabilistic wind power forecasting framework using deep CNNs and conformalized regression forests. Appl Energy 2024;361:122900. \url{https://doi.org/10.1016/j.apenergy.2024.122900}

\bibitem[22]{ref22}
Bi K, Xie L, Zhang H, Chen X, Gu X, Tian Q. Accurate medium-range global weather forecasting with 3D neural networks. Nature 2023;619:533-8. \url{https://doi.org/10.1038/s41586-023-06185-3}

\bibitem[23]{ref23}
Fan H, Li M, Zhang Z, Cheng L, Ye Y, Liu D. M2WLLM: multi-modal multi-task ultra-short-term wind power prediction algorithm based on large language model. arXiv preprint arXiv:2506.00531; 2025.

\bibitem[24]{ref24}
Zhou H, Zhang S, Peng J, Zhang S, Li J, Xiong H, et al. Informer: beyond efficient transformer for long sequence time-series forecasting. Proc AAAI Conf Artif Intell 2021;35:11106-15.

\bibitem[25]{ref27}
Nie Y, Nguyen NH, Sinthong P, Kalagnanam J. A time series is worth 64 words: long-term forecasting with transformers. Int Conf Learn Represent; 2023.

\bibitem[26]{ref28}
Zhang Y, Yan J. Crossformer: transformer utilizing cross-dimension dependency for multivariate time series forecasting. Int Conf Learn Represent; 2023.

\bibitem[27]{ref34}
Liu Y, Wu H, Wang J, Long M. Non-stationary transformers: exploring the stationarity in time series forecasting. Adv Neural Inf Process Syst 2022;35:9881-93.

\bibitem[28]{ref37}
Wu X, Zhan H, Hu J, Wang Y. Non-stationary GNNCrossformer: transformer with graph information for non-stationary multivariate spatio-temporal wind power data forecasting. Appl Energy 2025;377:124492. \url{https://doi.org/10.1016/j.apenergy.2024.124492}

\bibitem[29]{ref38}
Hou G, Li Q, Huang C. Spatiotemporal forecasting using multi-graph neural network assisted dual-domain transformer for wind power. Energy Convers Manag 2025;325:119393. \url{https://doi.org/10.1016/j.enconman.2024.119393}

\bibitem[30]{ref39}
Lin S, Wang S, Xu X, Li R, Shi P. GAOformer: an adaptive spatiotemporal feature fusion transformer utilizing GAT and optimizable graph matrixes for offshore wind speed prediction. Energy 2024;292:130404. \url{https://doi.org/10.1016/j.energy.2024.130404}

\bibitem[31]{ref40}
Chen J, Fu X, Zhang L, Shen H, Wu J. A novel offshore wind power prediction model based on TCN-DANet-sparse transformer and considering spatio-temporal coupling in multiple wind farms. Energy 2024;308:132899. \url{https://doi.org/10.1016/j.energy.2024.132899}

\bibitem[32]{ref41}
Lim B, Arık SÖ, Loeff N, Pfister T. Temporal fusion transformers for interpretable multi-horizon time series forecasting. Int J Forecast 2021;37:1748-64. \url{https://doi.org/10.1016/j.ijforecast.2021.03.012}

\bibitem[33]{ref42}
Niu Z, Han X, Zhang D, Wu Y, Lan S. Interpretable wind power forecasting combining seasonal-trend representations learning with temporal fusion transformers architecture. Energy 2024;306:132482. \url{https://doi.org/10.1016/j.energy.2024.132482}

\bibitem[34]{ref43}
Zhang C, Xiao X, Wang Y, Hou MZ, Huang S, Hu W, et al. Fine-grained ultra-short-term wind power forecasting based on Temporal Fusion Transformers integrated with turbine power time series clustering. Energy 2025;335:137995. \url{https://doi.org/10.1016/j.energy.2025.137995}

\bibitem[35]{ref35}
Carion N, Massa F, Synnaeve G, Usunier N, Kirillov A, Zagoruyko S. End-to-end object detection with transformers. In: European Conference on Computer Vision; 2020, p. 213-29. \url{https://doi.org/10.1007/978-3-030-58452-8_13}

\bibitem[36]{ref65}
Jaegle A, Gimeno F, Brock A, Vinyals O, Zisserman A, Carreira J. Perceiver: general perception with iterative attention. Proc Int Conf Mach Learn, PMLR 139; 2021, p. 4651-64.

\bibitem[37]{ref36}
Jaegle A, Borgeaud S, Alayrac JB, Doersch C, Ionescu C, Ding D, et al. Perceiver IO: a general architecture for structured inputs \& outputs. Int Conf Learn Represent; 2022.

\bibitem[38]{ref60}
Weston J, Chopra S, Bordes A. Memory networks. Int Conf Learn Represent; 2015.

\bibitem[39]{ref61}
Sukhbaatar S, Szlam A, Weston J, Fergus R. End-to-end memory networks. Adv Neural Inf Process Syst 2015;28:2440-8.

\bibitem[40]{ref66}
Alayrac JB, Donahue J, Luc P, Miech A, Barr I, Hasson Y, et al. Flamingo: a visual language model for few-shot learning. Adv Neural Inf Process Syst 2022;35:23716-36.

\bibitem[41]{ref70}
Perez E, Strub F, De Vries H, Dumoulin V, Courville A. FiLM: visual reasoning with a general conditioning layer. Proc AAAI Conf Artif Intell 2018;32.

\bibitem[42]{ref72}
Xu W, Liu Y, Fan X, Shen Z, Wu Q. Short-term wind power forecasting based on dual attention mechanism and gated recurrent unit neural network. Front Energy Res 2024;12:1346000. \url{https://doi.org/10.3389/fenrg.2024.1346000}

\bibitem[43]{ref73}
Naveed M, Ullah K, Shah SZ, Akram W, Ullah Z, Khan B. Enhancing wind power forecasting in power grids with dual attention-based LSTM mechanism. IET Gener Transm Distrib 2025;19:e70134. \url{https://doi.org/10.1049/gtd2.70134}

\end{thebibliography}


\end{document}
